% SPDX-License-Identifier: MIT
\chapter{Kurznachschlagewerk}\label{chap:appendix}
\ChapterSummary{Zum schnellen Nachschlagen finden Sie
hier die wichtigsten Einstiegspunkte, Build-Befehle,
Header-Dateien und eine kompakte Begriffsliste.
Die Aufstellung ist bewusst keine unveraenderliche,
vollständige API-Spezifikation.}

\section{Kompakte Orientierungsübersicht}

\begin{tabularx}{\linewidth}{@{}p{39mm}X@{}}
\toprule
Sie suchen & Beginnen Sie hier\\\midrule
Terminal-Beispiel & \filepath{examples/terminal_form_demo.cpp}\\
Grafik-Beispiel & \filepath{examples/rendered_sdl3_form_demo.cpp}\\
Core-API & \filepath{include/sasd/ui/}\\
Terminal-API & \filepath{include/sasd/ui/terminal/}\\
Rendered-API & \filepath{include/sasd/ui/rendered/}\\
CMake-Paket & \filepath{CMakeLists.txt} und \filepath{cmake/}\\
Entscheidungen & \filepath{docs/adr/README.md}\\
Roadmap & \filepath{docs/de/ROADMAP.md}\\
Tests & \filepath{tests/}\\
\bottomrule
\end{tabularx}

\section{Übliche Build-Kommandos}
\begin{lstlisting}[language=bash]
cmake -S . -B build -DSASD_UI_BUILD_TESTS=ON
cmake --build build --parallel
ctest --test-dir build --output-on-failure
\end{lstlisting}
Fuer Debug/Release/Visual-Studio-Besonderheiten
siehe Kapitel \ref{chap:install}.

\section{Ausgewählte Klassen und Header}

\begin{tabularx}{\linewidth}{@{}p{42mm}X@{}}
\toprule
Klasse / Thema & Header relativ zu include/\\\midrule
Window / Container & \filepath{sasd/ui/window.hpp}, \filepath{sasd/ui/container.hpp}\\
Label / Button / TextField & \filepath{sasd/ui/label.hpp},
 \filepath{sasd/ui/button.hpp}, \filepath{sasd/ui/text_field.hpp}\\
Checkbox / Radio & \filepath{sasd/ui/check_box.hpp}, \filepath{sasd/ui/radio_button.hpp}\\
ComboBox & \filepath{sasd/ui/combo_box.hpp}\\
VBox / HBox & \filepath{sasd/ui/vbox.hpp}, \filepath{sasd/ui/hbox.hpp}\\
Weitere Layouts & \filepath{sasd/ui/grid_layout.hpp}, \filepath{sasd/ui/form_layout.hpp}\\
Command / Shortcut & \filepath{sasd/ui/command.hpp}, \filepath{sasd/ui/shortcut.hpp}\\
Menues & \filepath{sasd/ui/menu_model.hpp}\\
Listen / Tabellen & \filepath{sasd/ui/list_model.hpp}, \filepath{sasd/ui/table_model.hpp}\\
Baum & \filepath{sasd/ui/tree_model.hpp}, \filepath{sasd/ui/tree_view.hpp}\\
\bottomrule
\end{tabularx}

\section{Begriffsliste}
\begin{description}[style=nextline]
\item[Backend] Bindeglied zu einem konkreten
Eingabe- und Ausgabesystem.
\item[Component] Objekt mit ausdruecklicher
Besitz- und Lebensdauersemantik.
\item[Container] Besitzt Komponenten; sichtbare
Widgets bilden einen untergeordneten Baum.
\item[Command] Nichtvisuelle Aktion, die
von mehreren UI-Ausloesern genutzt werden kann.
\item[Event] Eingangsereignis, etwa Key,
TextInput, Pointer oder Resize.
\item[MeasurementContext] Liefert Metriken
einer konkreten Praesentationsumgebung.
\item[Model] Stellt Daten bereit und meldet Aenderungen.
\item[SelectionModel] Haelt die aktuelle Auswahl
getrennt von der Darstellung.
\item[PresentationSink] Erzeugt darstellbare
Information aus Widget-Zustaenden.
\item[Viewport] Sichtbarer Ausschnitt einer
groesseren Datenmenge.
\item[Widget] Sichtbares semantisches
Bedienelement im Core.
\end{description}

\section{Hinweise zur Pflege des Handbuchs}
Neue oder umbenannte APIs muessen gegen
Header und Tests abgeglichen werden.
Fuer jede Aenderung an dieser deutschen
Ausgabe kann eine Sprachversion ein
entsprechendes Kapitel unter
\filepath{docs/manual/content/<lang>/}
erhalten. Ausstehende Übersetzungen
sollen sichtbar bleiben; eine
automatische Rueckfall-Uebersetzung
ist kein Ersatz für fachliche
Pruefung. Ein Kapitel darf nicht
eine M5-Planung als bereits
lieferfaehige M2-Funktion ausgeben.
